%----------------------------------
\section{Data Set}
\label{sec:Data Set}
%----------------------------------

\subsection{Public Data Sets and Attributes}

\textbf{Dataset 1 --- UCI Online Retail:} 541,909 invoice lines from a UK non-store retailer (1~December~2010--9~December~2011), with eight recorded columns: invoice, product, description, quantity, date, unit price, customer ID and country. \textbf{Dataset 2 --- UCI Online Shoppers Purchasing Intention:} 12,330 web sessions with 17 predictors plus the Revenue target. Predictors capture page counts and durations, bounce and exit rates, page value, special-day proximity, month, operating system, browser, region, traffic source, visitor type and weekend. Both are publicly accessible and licensed CC BY~4.0~\cite{chen2015onlineretail,sakar2018onlineshoppers}.

\begin{itemize}
    \item \href{https://archive.ics.uci.edu/dataset/352/online+retail}{UCI Online Retail --- accessible dataset page}
    \item \href{https://archive.ics.uci.edu/dataset/468/online+shoppers+purchasing+intention+dataset}{UCI Online Shoppers Purchasing Intention --- accessible dataset page}
\end{itemize}

\wordcount{88}

\subsection{Selection Rationale}

They collectively complement the insight into customers, forecasting, and dashboarding aspects of the role across two time horizons. Retail transactions provide insights into future customers by predicting who will buy again, which helps in planning resource allocation. Session data provides insights into the conversion of visitors into buyers, highlighting any friction in the digital journey and visits with higher intent to buy. They are not combined as they pertain to different firms and departments.

\wordcount{74}
